\documentclass[10pt,a4paper]{amsart}
\usepackage[margin=19mm]{geometry}
\usepackage{amsmath,amssymb,mathtools}
\usepackage[scr=rsfs,cal=euler]{mathalfa}
\usepackage{xfrac}
\usepackage[unicode,hidelinks,bookmarks=false]{hyperref}
\newcommand{\ie}{i.e.\ }
\newcommand{\cf}{cf.\ }
\newcommand{\resp}{respectively\ }
\newcommand{\Subsection}{\subsection}
\providecommand{\nobreakdash}{\nobreak\hbox{-}}
\setlength{\emergencystretch}{2em}
\multlinegap0pt
\usepackage{mathtools}
\usepackage{bm}
\usepackage{amssymb}
\usepackage{algorithm}\usepackage{algpseudocode}
\newtheorem{theorem}{Theorem}
\def\d{\mbox{d}}
\title{An adaptive radial basis function approach for efficiently solving multidimensional spatiotemporal integrodifferential equations}
\author{Mingtao Xia, Qijing Shen}
\date{Source: arXiv:2604.05276v1 (CC BY 4.0)}
\begin{document}
\maketitle

\noindent\emph{Source and licence.} An adaptive radial basis function approach for efficiently solving multidimensional spatiotemporal integrodifferential equations. Version arXiv:2604.05276v1 (CC BY 4.0), distributed by arXiv under the Creative Commons Attribution 4.0 International licence, http://creativecommons.org/licenses/by/4.0/, as recorded in the arXiv licence metadata for this exact version and shown on the abstract page https://arxiv.org/abs/2604.05276v1. Full licence notice: \texttt{licenses/arxiv-2604.05276-CC-BY-4.0.txt}.\par\smallskip

\noindent\textit{Editorial scope.} Counted unit: the article's theorem theorem2 (source0.tex lines 272--290). The article's complete original proof of it is delivered with it (source0.tex lines 292--311, one \texttt{proof} environment of 20 lines). No mathematics is written, repaired or re-derived: the statement and the proof are the article's own lines, in the article's own notation, and no retained line is substituted or re-flowed.  Retained uncounted as the source-local context the proof needs: lines 139--159.  Nothing was omitted from any retained span, so the package carries no omission record.  The retained lines cite 1 works, whose entries are transcribed verbatim from the article's own bibliography source in the e-print and delivered with short editorial labels recorded in the item's reference mapping.  No retained line contains a figure environment, an image-inclusion command or drawing code, so no image is delivered and the package declares no asset dependency.  Editorial formatting only, in addition: an AMS \texttt{amsart} preamble that restates the article's own declarations and macros used by the retained lines, namely the environments \texttt{theorem} on separate counters, together with the standard packages those lines need.  The retained environments print in this document as Theorem 1 in order of appearance, whereas the article prints the same items with the numbering of its own full document.  The complete native source of the article as distributed in the arXiv e-print for this exact version is bundled unchanged as \texttt{sources/197946/source0.tex}.
\section{Adaptive RBF approach for solving multidimensional spatiotemporal integrodifferential equations}
\label{section2}
In this section, we present and analyze the adaptive strategies used to apply radial basis functions for solving multidimensional spatiotemporal equations. For simplicity, in our theoretical analysis, we focus on the following general form of a spatiotemporal equation subject to homogeneous Dirichlet boundary conditions:
\begin{equation}
\begin{aligned}
    &\frac{\partial u(\bm{x}, t)}{\partial t} = A(u(\bm{x}, t), \bm{x}, t), \bm{x}, t), \,\,\bm{x}\in\Omega, t\in[0, T],\\
    & u(\bm{x}, 0) = u_0(\bm{x}), \,\, u(\bm{x}, t) = 0, \,\,\bm{x}\in\partial\Omega.
    \end{aligned}
    \label{model_problem}
\end{equation}
In Eq.~\eqref{model_problem}, $A$ denotes a general integrodifferential operator acting only on the spatial variable $\bm{x}$. We assume that the model problem \eqref{model_problem} is well posed and admits a unique solution $u(\bm{x},t)\in C^1([0,T],F_k^d)$ for some $k\ge 1$, where $F_k^d$ denotes the function space introduced in \cite{barthelmann2000high}:
\begin{equation}
\begin{aligned}
&F_k^d \coloneqq  \Big\{ f \coloneqq [-1, 1]^d \to \mathbb{R}  \big|\ D^{\bm{\alpha}}f \text{ continuous if } \alpha_i \leq k \text{ for all } i,\\
&\hspace{5cm}\bm{\alpha}\coloneqq(\alpha_1,...,\alpha_d)\in \mathbb{N}_0^d \Big\}
\end{aligned}
\end{equation}
equipped with the norm:
\begin{equation}
\|f\|_{k, \infty} = \max\left\{ \left\|D^{\alpha}u\right\|_{\infty} \ \middle|\ \alpha \in \mathbb{N}_0^d,\ \alpha_i \leq k \right\}.
\end{equation}

\begin{theorem}
\label{theorem2}
\rm
Assume $A$ in Eq.~\eqref{model_problem} is a linear operator and satisfies the following coercivity and linearity conditions: 
\begin{equation}
\begin{aligned}
    \int_{\Omega} u(\bm{x}, t)A(u(\bm{x}, t), \bm{x}, t)\d\bm{x}\leq  L\|u\|^2,\,\, A(u+v, \bm{x}, t) = A(u, \bm{x}, t)+ A(v, \bm{x}, t),
    \end{aligned}
    \label{A_condition}
\end{equation}
where $G$ is a function on $u$ as well as its derivatives on the boundary $\partial\Omega$.  
Then, the following error bound holds:
\begin{equation}
\begin{aligned}
    \big\|u(\bm{x}, t)-u_N(\bm{x}, t)\big\|^2&\leq \exp\big((2L+1)t\big)\cdot \bigg(\int_0^t \|A(u_N,\bm{x}, s)-\hat{A}(u_N, \bm{x}, s)\|^2\d s \\
    &\hspace{4cm} + \|u(\bm{x}, 0) - u_N(\bm{x}, 0)\|^2\bigg).
    \end{aligned}
\end{equation}
\end{theorem}

\begin{proof}
    We denote $f(t)\coloneqq\|u(\bm{x}, t) - u_N(\bm{x}, t)\|^2$. We have:
    \begin{equation}
    \begin{aligned}
        \partial_t f(t) &= 2(A(u, \bm{x},t) - A(u_N, \bm{x},t), u-u_N) + 2(A(u_N, \bm{x},t) - \hat{A}(u_N, \bm{x},t), u - u_N)\\
        &\quad\leq 2L(u-u_N, u-u_N) +  2\|A(u_N, \bm{x}, t)-\hat{A}(u_N, \bm{x}, t)\|\cdot\|u(\bm{x}, t) - u_N(\bm{x}, t)\|\\
        &\quad\quad\leq (2L+1)\|u-u_N\|^2 + \|A(u_N, \bm{x}, t)-\hat{A}(u_N, \bm{x}, t)\|^2
        \end{aligned}
        \label{model}
    \end{equation}
    In Eq.~\eqref{model}, $(u, v)\coloneqq \int_\Omega u\cdot v\d\bm{x}$. Applying the Gronwall's inequality to Eq.~\eqref{model}, we have:
    \begin{equation}
        \begin{aligned}
        \|u(\bm{x}, t)-u_N(\bm{x}, t)\|^2\leq \exp\big((2L+1)t\big)\cdot \bigg(2 \int_0^t 
         \|A(u_N, \bm{x}, s)-\hat{A}(u_N, \bm{x}, s)\|^2\d\bm{s} \\+ \|u(\bm{x}, 0) - u_N(\bm{x}, 0)\|^2\bigg),
        \end{aligned}
        \label{temporal_error_bound}
    \end{equation}
    which proves Theorem~\ref{theorem2}.
\end{proof}

\begin{thebibliography}{99}
\bibitem[barthe]{barthelmann2000high} Barthelmann, V., Novak, E., and Ritter, K., High dimensional polynomial interpolation on sparse grids, \textit{Advances in Computational Mathematics}, 12(4):273--288 (2000).
\end{thebibliography}
\end{document}
